\documentclass[a4paper,12pt]{article}
\usepackage[T2A]{fontenc}
\usepackage[utf8]{inputenc}
\usepackage[russian]{babel}
\usepackage{amsmath,amssymb,mathtools}
\renewcommand{\familydefault}{\sfdefault}
\usepackage{icomma}
\usepackage[a4paper,margin=19mm,headsep=7mm,footskip=12mm]{geometry}
\usepackage{microtype}
\usepackage{tikz}
\usetikzlibrary{arrows.meta,positioning,shapes.geometric,calc}
\usepackage{tabularx,booktabs,array}
\usepackage{enumitem}
\usepackage{hyperref}
\usepackage{fancyhdr}
\usepackage{setspace}
\usepackage{xcolor}

\definecolor{accent}{HTML}{2D6F73}
\definecolor{darkaccent}{HTML}{194A4D}
\definecolor{textgray}{HTML}{2D3338}
\definecolor{softgray}{HTML}{6C757D}
\definecolor{pale}{HTML}{EEF5F5}

\renewcommand{\today}{05.10.2026}
\hypersetup{hidelinks}
\setlength{\parindent}{0pt}
\setlength{\parskip}{0.42em}
\setlength{\headheight}{25pt}
\setlength{\emergencystretch}{2em}
\onehalfspacing
\raggedbottom
\color{textgray}

\pagestyle{fancy}
\fancyhf{}
\fancyhead[L]{\small\color{softgray}Чек-лист: задачи на производительность}
\fancyhead[R]{\small\color{softgray}05.10.2026}
\fancyfoot[C]{\small\color{softgray}\thepage}
\renewcommand{\headrulewidth}{0.3pt}
\renewcommand{\headrule}{\hbox to\headwidth{\color{accent}\leaders\hrule height \headrulewidth\hfill}}

\newcommand{\sectionrule}{\par\vspace{-0.25em}\noindent\textcolor{accent}{\rule{\linewidth}{0.6pt}}\par\vspace{0.25em}}
\newcommand{\key}[1]{\textcolor{darkaccent}{\textbf{#1}}}
\newcommand{\formula}[1]{\[\boxed{#1}\]}

\begin{document}

% ---------- TITLE ----------
\thispagestyle{empty}
\begin{center}
\vspace*{24mm}
{\Large\color{softgray}Математика}\par
\vspace{12mm}
{\Huge\bfseries\color{darkaccent}Чек-лист}\par
\vspace{5mm}
{\LARGE\bfseries Задачи на производительность\par}
\vspace{3mm}
{\large совместная работа, трубы, бригады и попарные условия}\par
\vspace{16mm}
{\large\color{softgray}\today}\par
\vfill
{\large\bfseries Лёвин Артём Александрович}\par
\vspace{2mm}
{\normalsize эксклюзивно для Анастасии Павловой}\par
\vspace{24mm}
\end{center}

\begin{tikzpicture}[remember picture,overlay]
\draw[accent,line width=1.25pt]
  ([xshift=18mm,yshift=16mm]current page.south west)
  .. controls ([xshift=72mm,yshift=31mm]current page.south west)
  and ([xshift=-68mm,yshift=7mm]current page.south east)
  .. ([xshift=-18mm,yshift=20mm]current page.south east);
\end{tikzpicture}
\clearpage

% ---------- CORE ----------
\section*{1. Базовая модель}
\sectionrule

В задачах на работу удобно держать перед глазами три величины:
\formula{A=p\,t}
где $A$ --- объём выполненной работы, $p$ --- производительность, $t$ --- время.

Если речь идёт об \key{одной целой работе} (один резервуар, один заказ, один забор), удобно принять $A=1$. Тогда
\[
p=\frac{1}{t},\qquad t=\frac{1}{p}.
\]

\begin{table}[h]
\centering
\renewcommand{\arraystretch}{1.3}
\begin{tabularx}{\linewidth}{@{}>{\raggedright\arraybackslash\bfseries}p{0.23\linewidth} X X@{}}
\toprule
Фраза в условии & Что это обычно означает & Куда записывать \\
\midrule
«за 5 часов», «за 40 минут» & длительность работы & в столбец $t$ \\
«12 деталей в час», «4 м$^3$ в минуту» & результат за единицу времени & в столбец $p$ \\
«весь резервуар», «один заказ» & общий объём работы & в столбец $A$; часто $A=1$ \\
«на 3 часа быстрее» & сравниваются времена & $x$ и $x-3$ \\
«на 4 детали в час больше» & сравниваются производительности & $x$ и $x+4$ \\
\bottomrule
\end{tabularx}
\end{table}

\key{Главная привычка.} Сначала определить, какая величина названа в условии: время, производительность или объём работы. Только после этого переносить число в таблицу.

\section*{2. Совместная работа и противоположные процессы}
\sectionrule

Если исполнители одновременно делают одну и ту же работу в одном направлении, их производительности складываются:
\[
p_{\text{совм}}=p_1+p_2+\dots
\]
Например, если первый выполняет работу за $a$ часов, а второй за $b$ часов, то
\[
\frac1a+\frac1b=\frac1{t_{\text{совм}}}.
\]

Если процессы направлены противоположно (например, одна труба наполняет бассейн, другая откачивает воду), производительности учитываются со знаками:
\[
p_{\text{итог}}=p_{\text{наполнение}}-p_{\text{откачивание}}.
\]
Знак определяется \key{смыслом процесса}, а не положением числа в тексте.

\section*{3. Ограничения и проверка смысла}
\sectionrule

Время и производительность в обычных задачах положительны. Поэтому при обозначении, например, времён $x$ и $x-8$ обязательно записывайте
\[
x>8.
\]
Это сразу помогает отбросить формально полученный, но бессмысленный корень.

После решения проверьте:
\begin{itemize}[leftmargin=1.5em,itemsep=0.2em,topsep=0.2em]
\item более производительный исполнитель должен тратить меньше времени на тот же объём работы;
\item при совместной работе в одном направлении итоговое время меньше времени каждого исполнителя по отдельности;
\item единицы времени и производительности должны быть согласованы;
\item найденный корень должен удовлетворять всем ограничениям исходной модели.
\end{itemize}

\clearpage

% ---------- EXAMPLES ----------
\section*{4. Типовые конструкции}
\sectionrule

\subsection*{4.1. Времена отличаются на известное число}
Пусть второй насос заполняет резервуар за $x$ часов, первый --- на $8$ часов быстрее, а вместе они заполняют резервуар за $7{,}5$ часа. Тогда
\[
\frac{1}{x-8}+\frac{1}{x}=\frac{1}{7{,}5}=\frac{2}{15},\qquad x>8.
\]
После приведения к общему знаменателю получаем корни $x=3$ и $x=20$. Ограничение $x>8$ оставляет $x=20$.

\key{Ответ по модели:} второй насос --- $20$ ч, первый --- $12$ ч.

\subsection*{4.2. Производительности отличаются на известное число}
Если один станок производит $x$ деталей в час, а другой --- $x+5$, то время на изготовление одинаковой партии из $N$ деталей равно
\[
\frac{N}{x}\quad\text{и}\quad\frac{N}{x+5}.
\]
Если более быстрый станок заканчивает на $\Delta t$ часов раньше, уравнение удобно писать в смысловом порядке:
\[
\frac{N}{x}-\frac{N}{x+5}=\Delta t.
\]
Сначала стоит мысленно проверить знак: у более медленного станка время действительно больше.

\subsection*{4.3. Изменился состав бригад}
Пусть рабочие имеют одинаковую квалификацию, а производительность одного рабочего равна $p$. Тогда производительность бригады из $n$ человек равна $np$.

В задаче, где две бригады выполняют \key{одинаковые заказы}, после изменения состава удобно разбить работу на два этапа: «до перехода» и «после перехода». Для разобранной модели с $16$ и $25$ рабочими, переходом $8$ человек через $7$ дней и одинаковым временем работы после перехода $y$:
\[
7\cdot16p+24py=7\cdot25p+17py.
\]
Поскольку $p>0$, общий множитель можно сократить. Получаем $y=9$, значит полный срок выполнения заказа равен $7+9=16$ дням.

\key{Смысл равенства:} левая и правая части --- полные объёмы двух одинаковых заказов.

\subsection*{4.4. Известны времена попарной работы трёх исполнителей}
Пусть индивидуальные производительности равны $x,y,z$. Если пары выполняют одну работу за $6$, $12$ и $8$ часов, то
\[
x+y=\frac16,\qquad y+z=\frac1{12},\qquad z+x=\frac18.
\]
Складывать эти равенства выгоднее, чем искать каждую производительность отдельно:
\[
2(x+y+z)=\frac16+\frac1{12}+\frac18=\frac38.
\]
Отсюда
\[
x+y+z=\frac3{16},\qquad t=\frac{1}{3/16}=\frac{16}{3}\text{ ч}=320\text{ мин}.
\]

\key{Идея:} ищите прежде всего ту величину, от которой непосредственно зависит ответ.

\section*{5. Типичные ошибки}
\sectionrule
\begin{enumerate}[leftmargin=1.7em,itemsep=0.25em,topsep=0.2em]
\item Прибавлять «4 м$^3$/ч» к времени. Это изменение \textbf{производительности}.
\item Складывать времена совместной работы. Складываются \textbf{производительности}.
\item Игнорировать направление процесса при заполнении/откачивании.
\item Смешивать часы и минуты в одном уравнении без перевода единиц.
\item Забывать, что одинаковая квалификация рабочих позволяет записать производительность бригады как $np$.
\item Делить на переменную величину без проверки, что она не равна нулю.
\item Принимать оба корня квадратного уравнения без проверки физического смысла.
\end{enumerate}

\clearpage

% ---------- FLOWCHART ----------
\thispagestyle{fancy}
\begin{center}
{\Large\bfseries\color{darkaccent}Алгоритм: задача на производительность}
\end{center}
\vspace{5mm}

\tikzset{
startstop/.style={ellipse,draw=accent,line width=0.9pt,align=center,minimum width=42mm,minimum height=9mm,fill=pale},
process/.style={rounded corners=2mm,rectangle,draw=accent,line width=0.9pt,align=center,text width=112mm,minimum height=10mm},
decision/.style={diamond,draw=accent,line width=0.9pt,align=center,text width=41mm,inner sep=1.3mm,aspect=2.2},
branch/.style={rounded corners=2mm,rectangle,draw=accent,line width=0.9pt,align=center,text width=54mm,minimum height=10mm,fill=pale},
arrow/.style={-{Stealth[length=3mm]},line width=0.9pt,color=darkaccent}
}

\begin{center}
\begin{tikzpicture}[node distance=5.5mm and 12mm]
\node (s) [startstop] {Прочитайте условие и определите искомую величину};
\node (u) [process,below=of s] {Разберите единицы: «за ... часов» $\to t$; «... в час» $\to p$; общий результат $\to A$};
\node (x) [process,below=of u] {Введите переменную и сразу запишите ограничения: время $>0$, знаменатели $\ne0$};
\node (tab) [process,below=of x] {Заполните таблицу «производительность --- время --- работа» по формуле $A=pt$};
\node (d) [decision,below=8mm of tab] {Процессы направлены одинаково?};
\node (plus) [branch,below left=6mm and 10mm of d] {Да: сложите производительности};
\node (minus) [branch,below right=6mm and 10mm of d] {Нет: задайте направление и вычтите противоположную производительность};
\node (eq) [process,below=12mm of d,yshift=-16mm] {Составьте уравнение по смыслу: равный объём, совместная производительность или сумма этапов};
\node (solve) [process,below=of eq] {Решите уравнение и отбросьте корни, нарушающие ограничения или смысл задачи};
\node (check) [process,below=of solve] {Проверьте единицы и здравый смысл результата; запишите ответ в требуемых единицах};
\node (f) [startstop,below=of check] {Готово};

\draw[arrow] (s)--(u);
\draw[arrow] (u)--(x);
\draw[arrow] (x)--(tab);
\draw[arrow] (tab)--(d);
\draw[arrow] (d)-- node[above left,font=\small]{да} (plus);
\draw[arrow] (d)-- node[above right,font=\small]{нет} (minus);
\draw[arrow] (plus.south) |- (eq.west);
\draw[arrow] (minus.south) |- (eq.east);
\draw[arrow] (eq)--(solve);
\draw[arrow] (solve)--(check);
\draw[arrow] (check)--(f);
\end{tikzpicture}
\end{center}
\clearpage

% ---------- TRAINING ----------
\section*{Тренировочный блок --- Путь А}
\sectionrule
{\color{softgray}10 задач по нарастающей сложности. Решения оформляйте через обозначение величин, таблицу или краткую модель и уравнение.}

\begin{enumerate}[leftmargin=1.8em,label=\textbf{\arabic*.},itemsep=0.75em]
\item Первый рабочий выполняет заказ за $12$ часов, второй --- за $18$ часов. За какое время они выполнят этот заказ, работая вместе?

\item Два насоса вместе заполняют резервуар за $6$ часов. Первый насос в одиночку заполняет его за $10$ часов. За сколько часов заполнит резервуар второй насос, работая один?

\item Один насос заполняет резервуар на $4$ часа быстрее другого. Вместе они заполняют резервуар за $4$ ч $48$ мин. За сколько часов каждый насос заполняет резервуар отдельно?

\item Одна труба заполняет бассейн за $4$ часа, а другая труба полностью откачивает воду из заполненного бассейна за $12$ часов. За сколько часов наполнится пустой бассейн, если открыть обе трубы одновременно?

\item Для изготовления партии из $120$ деталей первый станок делает на $5$ деталей в час больше второго и поэтому заканчивает работу на $2$ часа раньше. Найдите производительность каждого станка.

\item Две трубы вместе заполняют резервуар за $6$ часов. Первая труба в одиночку заполняет резервуар на $5$ часов быстрее второй. За сколько часов каждая труба заполнит резервуар отдельно?

\item Две бригады одинаково квалифицированных рабочих одновременно начали выполнять два одинаковых заказа. В первой бригаде было $14$ рабочих, во второй --- $22$. Через $5$ дней из второй бригады в первую перешли $6$ рабочих. После перехода обе бригады закончили свои заказы одновременно. Сколько дней потребовалось на выполнение каждого заказа?

\item Анна и Борис вместе выполняют работу за $8$ часов, Борис и Виктор --- за $12$ часов, Виктор и Анна --- за $8$ часов. За сколько часов они выполнят эту работу втроём?

\item Первый и второй мастера вместе выполняют работу за $6$ часов, второй и третий --- за $10$ часов, третий и первый --- за $7{,}5$ часа. За сколько часов выполнил бы эту работу каждый мастер отдельно?

\item Две трубы вместе заполняют резервуар за $4$ часа. Через $2$ часа совместной работы первую трубу закрыли, после чего вторая труба одна заполняла оставшуюся часть резервуара ещё $6$ часов. За сколько часов каждая труба заполняет резервуар в одиночку?
\end{enumerate}

\clearpage

% ---------- ANSWERS ----------
\section*{Ответы}
\sectionrule
\renewcommand{\arraystretch}{1.35}
\begin{table}[h]
\centering
\begin{tabularx}{\linewidth}{@{}>{\bfseries}p{0.08\linewidth} X@{}}
\toprule
№ & Ответ \\
\midrule
1 & $\dfrac{36}{5}$ ч $=7$ ч $12$ мин. \\
2 & $15$ ч. \\
3 & $8$ ч и $12$ ч. \\
4 & $6$ ч. \\
5 & $20$ деталей/ч и $15$ деталей/ч. \\
6 & $10$ ч и $15$ ч. \\
7 & $15$ дней. \\
8 & $6$ ч. \\
9 & $10$ ч, $15$ ч и $30$ ч. \\
10 & первая труба --- $6$ ч, вторая --- $12$ ч. \\
\bottomrule
\end{tabularx}
\end{table}

\vfill
\begin{center}
\small\color{softgray}
Лёвин Артём Александрович эксклюзивно для Анастасии Павловой
\end{center}

\end{document}
